\chapter{结论}
\label{cha:conclusion}

\section*{1. 论文工作总结}
\addcontentsline{toc}{section}{1. 论文工作总结}

本文围绕电梯轿厢中行人与电动车的双目标检测，探讨了 YOLO 类模型从训练环境迁移到边缘设备后的行为变化。研究的出发点是：离线训练精度并不能直接代表板端实际表现，模型格式转换、输出解释和后处理策略都会引入额外的不确定性。

针对这一问题，本文的主要工作可归纳为三个层面。在任务层面，将电梯安全视觉的核心需求——人员感知与电动车识别——统一为双目标检测框架，并系统梳理了轿厢场景中影响检测精度的遮挡、反光和近距离透视等因素。在迁移层面，将训练权重到板端推理的转换过程组织为五类约束检查，使每一步格式变换中的输入尺寸、类别顺序、坐标编码和阈值语义都可被逐项核实。在评价层面，不再依赖单一的离线 mAP，而是从分类别指标、阈值敏感性、NMS 参数影响、后处理消融、分段计时和视频连续输出等多角度交叉验证板端行为。

实验结果揭示了若干值得关注的现象。板端 mAP@0.5 相比训练阶段回退约 0.078，其中行人类别的下降幅度明显大于电动车类别（0.869 vs 0.952），这与轿厢近景遮挡导致行人候选框置信度偏低有关。阈值敏感性实验表明，总体 F1 在较宽的置信度区间内保持稳定，说明板端模型的安全输出不依赖精确的单一阈值调优。后处理消融则进一步揭示，应用层清理策略对电动车召回率的影响远大于模型本身——完整清理策略使电动车 recall 骤降至 0.765，这意味着部署一致性评价必须将后处理模块从模型权重效果中显式剥离。

总体而言，本文工作表明，边缘安全视觉系统的评价不应止步于“模型能跑、精度达标”，而应覆盖从格式转换到后处理策略的完整部署链路。本文整理的分阶段检查方法和多维评价思路，可为同类电梯或封闭空间的边缘视觉部署提供参考。

\section*{2. 工作展望}
\addcontentsline{toc}{section}{2. 工作展望}

围绕进一步应用，后续可从四个方向推进。第一，扩充电梯轿厢数据来源，覆盖更多楼宇类型、摄像机安装角度、光照条件和人车共现形态，使双目标检测模型在更丰富的真实场景中得到验证。第二，在更多 YOLO 变体和边缘芯片平台上复用本文的迁移检查流程，比较不同模型与硬件组合的部署行为，进一步分析方法在同类海思/昇腾边缘 AI 平台及工具链中的适用性。第三，开展训练、ONNX、OM 和板端四阶段的逐样本匹配，量化边框、类别分数和候选框数量在迁移过程中的变化，使部署一致性评价更加细粒度。第四，扩充带逐帧标注的视频评价数据，将连续输出稳定性与视频级检测精度结合起来，并在同一框架下继续拓展人数统计、跌倒检测等电梯安全视觉任务。
